\section*{Hidden Constants in Integration by Parts}

\url{https://www.youtube.com/watch?v=N5xtRr0jfWI} \\
\url{https://www.youtube.com/watch?v=V8uZYdd4Jnk}

\begin{itemize}

\item Why does every indefinite integral include an arbitrary constant $C$?

\item Derive the integration by parts formula from the product rule.

\item Why is the equation
  \[
    \int u\,dv = uv - \int v\,du
  \]
  technically incomplete if we ignore constants of integration? Where are the hidden constants in this equation?

\item Why can ignoring arbitrary constants sometimes lead to paradoxes or false identities when manipulating indefinite integrals?

\item Consider this example
  \[
    I = \int \underbrace{\frac{f'(x)}{f(x)}}_{uv'} \, dx =
      \underbrace{\frac{f(x)}{f(x)}}_{uv} - \int \underbrace{-\frac{f'(x)}{f^2(x)}}_{u'} \underbrace{f(x)}_{v} \, dx =
      1 + \int \frac{f'(x)}{f(x)} \, dx =
      1 + I
  \]
  \begin{align*}
    u  &= \frac{1}{f(x)}        & v' &= f'(x) \\
    u' &= -\frac{f'(x)}{f(x)^2} & v  &= f(x)
  \end{align*}
  Where is the error in this paradoxical derivation?

\item Consider this example
  \[
    I = \int f(x)g'(x) \, dx
  \]
  \begin{align*}
    u  &= f(x)  & v' &= g'(x) \\
    u' &= f'(x) & v  &= g(x) + A
  \end{align*}
  Show how all the constants of integration are absorbed into one when applying integration by parts.

\item Why can repeated applications of integration by parts accumulate multiple constants?

\item How can these multiple constants always be absorbed into a single constant in the final answer?

\item Why is it sometimes useful to deliberately choose a non-zero constant when computing an antiderivative inside integration by parts?

\item If the integration by part is applied naively to integral $\int x\arctan(x) \, dx$ it leads to an incomplete solution
  \[
    \int \underbrace{x}_{v'} \underbrace{\arctan(x)}_{u} \, dx
      = \underbrace{\frac{x^2}{2}}_{v} \underbrace{\arctan(x)}_{u} -
        \int \underbrace{\frac{x^2}{2}}_{v} \underbrace{\frac{1}{x^2+1}}_{u'}\,dx.
  \]
  How can choosing a convenient constant in the antiderivative simplify the naive solution?

\item If the integration by part is applied naively to integral $\int \ln(x+2) \, dx$ it leads to an incomplete solution
  \[
    \int \underbrace{1}_{v'} \underbrace{\ln(x+2)}_{u} \, dx
      = \underbrace{x}_{v} \underbrace{\ln(x+2)}_{u} -
        \int \underbrace{x}_{v} \underbrace{\frac{1}{x+2}}_{u'} \, dx.
  \]
  How can choosing a convenient constant in the antiderivative simplify the naive solution?

\item Why are definite integrals immune to the constant-of-integration ambiguity?

\end{itemize}
